Tato krátká sekce shrnuje zásady pro návrh výuky s generativní AI jako rámec před DPIA, smlouvami a technickým řešením.
Základní principy
\begin{itemize}
  \item Minimalizace dat; nezpracovávejte citlivá či přímo identifikovatelná data v externích službách \cite{kb_ai_101_tipu_pdf_04e4a115_page_8_1}.
  \item Účelové omezení, transparentnost a tam, kde možno, anonymizace/pseudonymizace.
  \item DPIA při plánování pilotů a širších nasazení; DPIA povinná při středních/vysokých rizicích \cite{kb_malinka_chatgpt_ve_vyuce_dva_roky_pote_2024_pdf_90bc3aaf_page_12_1}.
\end{itemize}
Kdy zapojit bezpečnost/PRIV: při zpracování osobních/citlivých dat, ukládání mimo EU nebo při pilotu sbírajícím behaviorální metriky či ovlivňující hodnocení; před spuštěním proveďte checklist (minimalizace PII, DPIA/odůvodnění, ověření DPA, lidská validace).